\documentclass[11pt]{article}
\usepackage[margin=1in]{geometry}
\usepackage[T1]{fontenc}
\usepackage[utf8]{inputenc}
\usepackage{amsmath,amssymb,amsthm}
\usepackage{enumitem}

\title{ICPC World Finals 2025\\L. Walking on Sunshine}
\author{}
\date{}

\begin{document}
\maketitle

\section*{Problem Summary}

We may move anywhere in the plane. Only movement with a southward component is painful, and even that
cost disappears inside the shaded rectangles. We must minimize the total painful distance needed to get
from the contest to the awards ceremony.

\section*{Key Observations}

\begin{itemize}[leftmargin=*]
    \item If the destination is not south of the start, the answer is immediately $0$: we can move only east,
    west, and north.
    \item Horizontal movement is always free. Northward movement is always free. Only the southward part
    of our motion matters.
    \item Because horizontal motion is free, whenever a shaded rectangle intersects some $y$-range we may
    move sideways for free and use that rectangle to descend through that entire vertical interval without
    paying.
    \item Therefore the $x$-coordinates are irrelevant. The answer is simply:
    the total required southward drop, minus the total length of the union of the shaded $y$-intervals that
    intersect that drop.
\end{itemize}

\section*{Algorithm}

\begin{enumerate}[leftmargin=*]
    \item If $y_a \ge y_c$, print $0$.
    \item Otherwise, the total unavoidable southward range is $[y_a,y_c]$.
    For each rectangle, intersect its vertical interval with this range.
    If the intersection is nonempty, store that interval.
    \item Sort the stored intervals by their lower endpoint and merge them.
    Let $U$ be the total merged length.
    \item The answer is
    \[
    (y_c-y_a)-U.
    \]
\end{enumerate}

\section*{Correctness Proof}

We prove that the algorithm returns the correct answer.

\paragraph{Lemma 1.}
Any path from the contest to the awards ceremony must have total southward displacement at least
$y_c-y_a$ when $y_a < y_c$.

\paragraph{Proof.}
The starting $y$-coordinate is $y_c$ and the final one is $y_a$.
Northward motion cannot reduce the required net drop, so the total southward component of the path is at
least $y_c-y_a$. \qed

\paragraph{Lemma 2.}
Every shaded $y$-interval intersecting $[y_a,y_c]$ can be traversed with zero pain.

\paragraph{Proof.}
Horizontal motion is free, so before descending through a shaded interval we may move sideways into the
corresponding rectangle at no cost. While inside shade, even southward movement is free. After leaving
that interval we may again move sideways for free. Thus the entire covered vertical interval contributes
zero painful distance. \qed

\paragraph{Theorem.}
The algorithm outputs the correct answer for every valid input.

\paragraph{Proof.}
By Lemma 1, any valid path must pay for at least the part of the necessary southward drop that is not
covered by shade. By Lemma 2, every shaded subinterval of $[y_a,y_c]$ can indeed be used to make that
part of the descent free. Therefore the minimum painful distance is exactly the uncovered length of
$[y_a,y_c]$, which is what the algorithm computes. \qed

\section*{Complexity Analysis}

Sorting the rectangle intervals costs $O(n \log n)$. Merging them is linear. The memory usage is $O(n)$.

\section*{Implementation Notes}

\begin{itemize}[leftmargin=*]
    \item Because the rectangles are disjoint, the merge step is especially simple, but the standard interval
    union code already handles the general case.
    \item The answer is printed as a floating-point number only because the judge allows that format; in this
    solution it is actually the length of a union of intervals.
\end{itemize}

\end{document}
